\setcounter{ex}{0}

% --- CÂU 1 ---
\begin{ex}
	Cho phương trình $2\cos^2 x - 3\cos x + 1 = 0\ (*)$. Các mệnh đề sau đúng hay sai?
	\choiceTFt
	{\True Phương trình $(*)$ tương đương với $\cos x = 1$ hoặc $\cos x = \dfrac{1}{2}$}
	{\True Tập nghiệm của $(*)$ là $\left\{k2\pi;\ \pm\dfrac{\pi}{3} + k2\pi,\ k \in \mathbb{Z}\right\}$}
	{\True Phương trình $(*)$ có đúng $4$ nghiệm thuộc đoạn $[0; 2\pi]$}
	{Tổng các nghiệm của $(*)$ thuộc đoạn $[0; 2\pi]$ bằng $3\pi$}
	\loigiai{
		\begin{itemchoice}
			\itemch \textbf{Đúng.} Ta có $2\cos^2 x - 3\cos x + 1 = 0 \iff (\cos x - 1)(2\cos x - 1) = 0 \iff \hoac{\cos x &= 1 \\ \cos x &= \dfrac{1}{2}.}$
			\itemch \textbf{Đúng.} Giải từng phương trình:
			$$\hoac{\cos x &= 1 \\ \cos x &= \dfrac{1}{2}} \iff \hoac{x &= k2\pi \\ x &= \pm\dfrac{\pi}{3} + k2\pi}\ (k \in \mathbb{Z}).$$
			\itemch \textbf{Đúng.} Trên đoạn $[0; 2\pi]$, phương trình có đúng $4$ nghiệm phân biệt là $0,\ \dfrac{\pi}{3},\ \dfrac{5\pi}{3},\ 2\pi$.
			\itemch \textbf{Sai.} Tổng các nghiệm thuộc đoạn $[0; 2\pi]$ bằng $0 + \dfrac{\pi}{3} + \dfrac{5\pi}{3} + 2\pi = 4\pi \ne 3\pi$.
		\end{itemchoice}
	}
\end{ex}

% --- CÂU 2 ---
\begin{ex}
	Cho hàm số $y = 3\cos\left(2x - \dfrac{\pi}{2}\right) + 1$. Các mệnh đề sau đúng hay sai?
	\choiceTFt
	{\True $y = 3\sin 2x + 1$}
	{Hàm số tuần hoàn với chu kì $2\pi$}
	{\True Giá trị lớn nhất của hàm số bằng $4$}
	{Hàm số đã cho là hàm số lẻ}
	\loigiai{
		\begin{itemchoice}
			\itemch \textbf{Đúng.} Vì $\cos\left(2x - \dfrac{\pi}{2}\right) = \cos\left(\dfrac{\pi}{2} - 2x\right) = \sin 2x$ nên $y = 3\sin 2x + 1$.
			\itemch \textbf{Sai.} Hàm số $y = 3\sin 2x + 1$ tuần hoàn với chu kì cơ sở là $T = \dfrac{2\pi}{2} = \pi \ne 2\pi$.
			\itemch \textbf{Đúng.} Vì $-1 \le \sin 2x \le 1$ nên $-2 \le 3\sin 2x + 1 \le 4$. Vậy giá trị lớn nhất của hàm số bằng $4$.
			\itemch \textbf{Sai.} Tại $x = 0$, ta có $y(0) = 1 \ne -y(-0) = -1$ nên hàm số đã cho không phải là hàm số lẻ.
		\end{itemchoice}
	}
\end{ex}

% --- CÂU 3 ---
\begin{ex}
	Cho cấp số cộng $(u_n)$ thỏa mãn $\heva{u_1 + u_3 &= -10 \\ u_2 + u_5 &= -19.}$ Các mệnh đề sau đúng hay sai?
	\choiceTFt
	{\True Công sai $d = -3$}
	{\True Số hạng đầu $u_1 = -2$}
	{Số hạng thứ mười $u_{10} = -28$}
	{\True Tổng $10$ số hạng đầu $S_{10} = -155$}
	\loigiai{
		\begin{itemchoice}
			\itemch \textbf{Đúng.} Gọi $d$ là công sai của cấp số cộng. Ta có hệ phương trình:
			$$\heva{u_1 + u_3 &= -10 \\ u_2 + u_5 &= -19} \iff \heva{2u_1 + 2d &= -10 \\ 2u_1 + 5d &= -19} \iff \heva{u_1 &= -2 \\ d &= -3.}$$
			Vậy công sai $d = -3$.
			\itemch \textbf{Đúng.} Số hạng đầu $u_1 = -2$.
			\itemch \textbf{Sai.} Số hạng thứ mười là $u_{10} = u_1 + 9d = -2 + 9 \cdot (-3) = -29 \ne -28$.
			\itemch \textbf{Đúng.} Tổng $10$ số hạng đầu là $S_{10} = 10u_1 + \dfrac{10 \cdot 9}{2}d = 10 \cdot (-2) + 45 \cdot (-3) = -155$.
		\end{itemchoice}
	}
\end{ex}

% --- CÂU 4 ---
\begin{ex}
	Cho hình chóp $S.ABCD$ có đáy là tứ giác lồi $ABCD$, $AB$ cắt $CD$ tại $F$, $AC$ cắt $BD$ tại $O$ (hình vẽ). Các mệnh đề sau đúng hay sai?
	\begin{center}
		\begin{tikzpicture}[scale=0.85, font=\footnotesize, line join=round, line cap=round]
			\coordinate (A) at (0,0);
			\coordinate (B) at (3.2,0);
			\coordinate (F) at (4.2,0);
			\coordinate (D) at (0.8,1.3);
			\coordinate (C) at ($(D)!0.65!(F)$);
			\coordinate (O) at (intersection of A--C and B--D);
			\coordinate (S) at (2.2,3.4);
			\draw[thick] (S)--(A) (S)--(B) (S)--(C) (S)--(F) (A)--(F) (B)--(C);
			\draw[dashed] (S)--(D) (A)--(D)--(F) (A)--(C) (B)--(D);
			\foreach \p/\pos in {S/above, A/below left, B/below, F/below right, C/above right, D/left, O/below}
				\node[\pos] at (\p) {$\p$};
			\foreach \p in {S,A,B,C,D,O,F}
				\fill (\p) circle (1.2pt);
		\end{tikzpicture}
	\end{center}
	\choiceTFt
	{$F \in (SBC)$}
	{\True Giao tuyến của $(SAB)$ và $(SCD)$ là đường thẳng $SF$}
	{\True $O$ thuộc giao tuyến của hai mặt phẳng $(SAC)$ và $(SBD)$}
	{Ba đường thẳng $AC, BD, SF$ đồng quy}
	\loigiai{
		\begin{itemchoice}
			\itemch \textbf{Sai.} Vì $(SBC) \cap (ABCD) = BC$, mà $F = AB \cap CD \notin BC$ nên $F \notin (SBC)$.
			\itemch \textbf{Đúng.} Ta có $S \in (SAB) \cap (SCD)$ và $F = AB \cap CD \in (SAB) \cap (SCD)$ nên $(SAB) \cap (SCD) = SF$.
			\itemch \textbf{Đúng.} Vì $O = AC \cap BD$ nên $O \in (SAC) \cap (SBD)$, tức là $O$ thuộc giao tuyến $SO$ của hai mặt phẳng $(SAC)$ và $(SBD)$.
			\itemch \textbf{Sai.} Hai đường thẳng $AC$ và $BD$ cắt nhau tại $O$ thuộc mặt phẳng $(ABCD)$, còn giao điểm của $SF$ với $(ABCD)$ là điểm $F \ne O$ nên ba đường thẳng $AC, BD, SF$ không đồng quy.
		\end{itemchoice}
	}
\end{ex}
